\clearpage
\SourceLeaf{95}%

\begin{center}
\large SCH\"ONES THEOREM DER WAHRSCHEINLICHKEITSRECHNUNG.
\end{center}

\noindent\hspace*{\parindent}Wenn [man setzt]:
\[
\int e^{itu}\varphi t\,dt = \psi u \mathord{.} \sqrt{2\pi},
\]
\noindent das Integral von $t=-\infty$ bis $t=+\infty$ ausgedehnt, so ist
\[
\int e^{-itu}\psi u\,du = \varphi t \mathord{.} \sqrt{2\pi},
\]
\noindent das Integral ebenso genommen.

Die Begr\"undung dieses Satzes ist in der Formel enthalten:
\[
\Sigma\,\varphi\!\left(t+\frac{2\pi k}{\omega}\right) = \frac{\omega}{2\pi}\Sigma\, e^{-ikt\omega}\int \varphi\theta\mathord{.} d\theta\mathord{.} e^{ik\theta\omega},
\]
\noindent wo das Integralzeichen die Ausdehnung $\theta=-\infty$ bis $\theta=+\infty$ voraussetzt, und das $\Sigma$-Zeichen sich auf alle ganzen positiven und negativen Werthe von $k$ bezieht, indem man darin nur $\omega$ unendlich abnehmen zu lassen braucht.

\Needspace{14\baselineskip}
\begin{center}
BEMERKUNGEN ZUM THEOREM DER WAHRSCHEINLICHKEITSRECHNUNG.
\end{center}

\noindent\hspace*{\parindent}Das vorliegende Theorem hat \textsc{Gauss} in den Einbanddeckel eines Buches eingetragen, welches den Titel »Opuscula mathematica, 1799---1813« tr\"agt. Eine genaue Zeitangabe, wann der Satz gefunden wurde, ist nicht vorhanden.

Man kann dem Theorem durch leichte Umrechnung auch die v\"ollig gleichwerthige Gestalt verleihen:
\[
2\pi\varphi(t) = \int^{+\infty}_{-\infty}\!\left(du\int^{+\infty}_{-\infty}\!\cos u(t-x)\mathord{.}\varphi(x)\,dx\right).
\]
